\begin{abstract}
A lower bias score on ambiguous items may reflect more abstention, not a
change in preferences. Using an exact identity implied by BBQ's metric, we
decompose 32 published method--model results: a median 80.4\% of the reported
improvement is accounted for by the coverage component (65.3\% under a
symmetric split), committed-answer bias \emph{worsened} in 7 pairs, and the
identified set widened in 31. First-party, on six checkpoints and eleven BBQ
categories, abstention varies mainly by checkpoint (83\% of variance) and the
stereotype-aligned rate by category (68\%), so the score ranks checkpoints
largely by abstention ($\rho=-0.83$, exact $p=0.058$), surviving reworded
prompts and extending to twelve checkpoints ($\rho=-0.92$) and six languages,
with one exception, Catalan, consistent with the identity. Abstainers still
lean toward the stereotype (56--64\%). The natural remedy, forced choice
checked for agreement, is unsound in its plug-in form: under a resampled null
it falsely rejects about 21\% of items at 16 draws, unimproved through
4{,}096. An exact-binomial correction gives per-item false-positive rates of
at most 0.05\% (none observed family-wise). Applied to 138{,}240 generations,
the certificate detects incompatibilities, which we classify by dissenting
condition as consistent with instrument failure or response-policy effects.
For Qwen2.5, the recovery gap strengthens under an alternative instrument that
shows no instrument-failure pattern.
\end{abstract}
